% Unit 5 self-study -- "I do / You do" ladder for CED 5.1-5.11.
% Twelve ladders, forty-eight questions.
%
% Disjoint from u05-frq, which uses the 2NO + O2 initial-rates table, the
% N2O5 three-step mechanism, a k = 1.5e-3 half-life and the collision-model
% temperature question. This document uses an A + 2B table, the NO2/CO
% mechanism, k = 0.0300 and different catalysis items.
%
% CED 5.6 excludes Arrhenius CALCULATIONS. Ladder 8 reads an energy
% profile and reasons qualitatively; no question extracts Ea from data.
\documentclass{apchem-worksheet}

\apcunit{5}
\apcunittitle{Kinetics}
\apcdoctitle{Self-Study \textbullet\ CED 5.1--5.11, I~do / You~do}
\apcsubtitle{Twelve ladders \textbullet\ four YOUR TURN questions each \textbullet\ work alone, check after all four}

\begin{document}
\apctitleblock

\begin{apcnote}[How to use these notes --- read this first]
Each skill is a \textbf{ladder}: a worked example, then \textbf{four}
YOUR TURN questions --- same skill, new numbers, no help. Work all four
before checking anything.

\textbf{The one idea this whole unit turns on.} Kinetics is about
\emph{how fast}, thermodynamics is about \emph{how far}. They are
independent. A reaction can be enormously favourable and still take
centuries; diamond turning to graphite is the standard example. Never
answer a rate question with an argument about stability.

\textbf{The habit that decides this unit.} Reaction orders are
\textbf{experimental}. They come from data, or from the slow step of a
mechanism --- never from the coefficients of the overall equation. Reading
orders off a balanced equation is the single most common way to lose marks
here.

\textbf{One scope note.} The CED excludes \textbf{Arrhenius equation
calculations}: you will not be asked to extract an activation energy from
two rate constants or from a plot of $\ln k$ against $1/T$. Reading an
energy profile and reasoning about temperature and catalysts qualitatively
both remain fully examinable.
\end{apcnote}

% =====================================================================
\section*{\sffamily Ladder 1 \textbullet\ Expressing reaction rate}
\ced{5.1}

Rate is a change in concentration per unit time. Divide each species'
rate by its coefficient and every species gives the \emph{same} number ---
one rate for the reaction.

\begin{workedexample}[I do: relating species]
For \ce{2N2O5 -> 4NO2 + O2}, oxygen appears at
\SI{0.0125}{\Molar\per\second}. Find the rates for the other two.

\textbf{Scale by the coefficients.} Relative to \ce{O2} (coefficient 1):

\ce{NO2} has coefficient 4, so it forms at
$4 \times 0.0125 = \mathbf{0.0500}$~M/s.

\ce{N2O5} has coefficient 2, so it is \emph{consumed} at
$2 \times 0.0125 = \mathbf{0.0250}$~M/s. As a rate of change,
$\Delta[\ce{N2O5}]/\Delta t = -0.0250$~M/s --- negative, because it is
disappearing.

\textbf{The single reaction rate:}
\[ -\tfrac{1}{2}\frac{\Delta[\ce{N2O5}]}{\Delta t}
   = \tfrac{1}{4}\frac{\Delta[\ce{NO2}]}{\Delta t}
   = \frac{\Delta[\ce{O2}]}{\Delta t} = \SI{0.0125}{\Molar\per\second} \]
\end{workedexample}

\begin{yourturn}[YOUR TURN 1 --- four questions]
\begin{enumerate}[leftmargin=*,itemsep=0.5em]
  \item For \ce{N2 + 3H2 -> 2NH3}, if \ce{N2} is consumed at
        \SI{0.020}{\Molar\per\second}, how fast is \ce{H2} consumed?
        \workspace[1.4cm]
  \item How fast is \ce{NH3} formed in (a)? \workspace[1.4cm]
  \item Why is the rate of change of a reactant negative?
        \workspace[1.4cm]
  \item Why divide each rate by its coefficient? \workspace[1.6cm]
\end{enumerate}
\selfcheck{(a) \SI{0.060}{\Molar\per\second} \quad
(b) \SI{0.040}{\Molar\per\second} \quad (c) its concentration falls \quad
(d) to get one rate independent of the species watched}
\keyonly{\begin{apcanswer}(a) $3 \times 0.020 = \mathbf{0.060}$~M/s.
(b) $2 \times 0.020 = \mathbf{0.040}$~M/s.
(c) Because its concentration is \emph{decreasing}, so
$\Delta[\text{reactant}]$ is negative. The reaction rate itself is quoted
as a positive number, which is why the expression carries a minus sign.
(d) Because the species are consumed and produced in fixed proportions set
by the coefficients. Dividing those proportions out leaves a single number
describing the reaction as a whole, the same whichever species was
monitored.\end{apcanswer}}
\end{yourturn}

% =====================================================================
\section*{\sffamily Ladder 2 \textbullet\ Rate law from initial rates}
\ced{5.2}

Compare two experiments in which \emph{one} concentration changes and the
others are held constant. The factor the rate changes by tells you the
order.

\begin{workedexample}[I do: reading a rate table]
For \ce{A + 2B -> C}:

\begin{center}
\begin{tabular}{cccc}
\hline
\textbf{Exp} & $[\ce{A}]_0$ & $[\ce{B}]_0$ & \textbf{rate} (M/s) \\
\hline
1 & 0.050 & 0.10 & $1.2\times10^{-3}$ \\
2 & 0.100 & 0.10 & $4.8\times10^{-3}$ \\
3 & 0.050 & 0.20 & $2.4\times10^{-3}$ \\
\hline
\end{tabular}
\end{center}

\textbf{Order in A:} compare 1 and 2, where $[\ce{B}]$ is constant.
$[\ce{A}]$ doubles and the rate goes up by
$4.8/1.2 = 4$. Since $2^n = 4$, $n = \mathbf{2}$.

\textbf{Order in B:} compare 1 and 3, where $[\ce{A}]$ is constant.
$[\ce{B}]$ doubles and the rate doubles, so $n = \mathbf{1}$.

\[ \text{rate} = k[\ce{A}]^2[\ce{B}] \qquad
   \text{overall order } 2 + 1 = \mathbf{3} \]

\textbf{Note what the equation would have suggested:} a coefficient of 2
on B, hence perhaps second order in B. The data say first. Orders are
experimental.
\end{workedexample}

\begin{yourturn}[YOUR TURN 2 --- four questions]
\begin{enumerate}[leftmargin=*,itemsep=0.5em]
  \item If tripling a reactant's concentration multiplies the rate by 9,
        what is the order? \workspace[1.4cm]
  \item If doubling a concentration leaves the rate unchanged, what is the
        order? \workspace[1.4cm]
  \item If doubling a concentration multiplies the rate by 8, what is the
        order? \workspace[1.4cm]
  \item Why must the other concentrations be held constant?
        \workspace[1.6cm]
\end{enumerate}
\selfcheck{(a) second \quad (b) zero \quad (c) third \quad (d) otherwise
you cannot attribute the change to one reactant}
\keyonly{\begin{apcanswer}(a) $3^n = 9$, so $n = \mathbf{2}$.
(b) $2^n = 1$, so $n = \mathbf{0}$ --- the rate does not depend on that
reactant at all.
(c) $2^n = 8$, so $n = \mathbf{3}$.
(d) Because the rate depends on all of them at once. If two
concentrations changed between the experiments, the change in rate could
not be attributed to either one alone --- the comparison would be
uninterpretable.\end{apcanswer}}
\end{yourturn}

% =====================================================================
\section*{\sffamily Ladder 3 \textbullet\ The rate constant and its units}
\ced{5.2}

$k$ is found by substituting one experiment back into the rate law. Its
\emph{units} are whatever makes the equation dimensionally consistent, and
they depend on the overall order.

\begin{workedexample}[I do: value and units together]
From experiment 1 above, find $k$ with units.

\[ k = \frac{\text{rate}}{[\ce{A}]^2[\ce{B}]}
     = \frac{1.2\times10^{-3}}{(0.050)^2(0.10)}
     = \frac{1.2\times10^{-3}}{2.5\times10^{-4}} = \mathbf{4.8} \]

\textbf{Units:} rate is \si{\Molar\per\second} and
$[\ce{A}]^2[\ce{B}]$ is \si{\Molar\cubed}, so $k$ carries
\si{\per\Molar\squared\per\second} --- that is,
M\textsuperscript{-2}\,s\textsuperscript{-1}.

\textbf{Check with another experiment:}
$4.8(0.100)^2(0.10) = 4.8\times10^{-3}$, matching experiment 2.
\checkmark A rate constant that fails this check means the orders are
wrong.

\textbf{The pattern:} overall order 0 gives \si{\Molar\per\second}; order
1 gives \si{\per\second}; order 2 gives
\si{\per\Molar\per\second}; order 3 gives
\si{\per\Molar\squared\per\second}.
\end{workedexample}

\begin{yourturn}[YOUR TURN 3 --- four questions]
\begin{enumerate}[leftmargin=*,itemsep=0.5em]
  \item Units of $k$ for a first-order reaction: \workspace[1.4cm]
  \item Units of $k$ for a second-order reaction: \workspace[1.4cm]
  \item A reaction has rate $= k[\ce{A}][\ce{B}]$ with rate
        \SI{6.0e-4}{\Molar\per\second} when both are
        \SI{0.20}{\Molar}. Find $k$. \workspace[1.8cm]
  \item Why do the units of $k$ differ between orders?
        \workspace[1.6cm]
\end{enumerate}
\selfcheck{(a) \si{\per\second} \quad (b) \si{\per\Molar\per\second} \quad
(c) \SI{1.5e-2}{\per\Molar\per\second} \quad (d) they must make the
equation dimensionally consistent}
\keyonly{\begin{apcanswer}(a) \textbf{\si{\per\second}}
(s\textsuperscript{-1}).
(b) \textbf{\si{\per\Molar\per\second}}
(M\textsuperscript{-1}\,s\textsuperscript{-1}).
(c) $k = 6.0\times10^{-4}/[(0.20)(0.20)] =
\mathbf{1.5\times10^{-2}}$~M\textsuperscript{-1}\,s\textsuperscript{-1}.
(d) Because rate always has units of \si{\Molar\per\second}, while the
concentration term on the right is raised to the overall order. $k$ must
supply whatever units are needed to reconcile the two, so its units change
with that order --- which is why quoting $k$ without units is
meaningless.\end{apcanswer}}
\end{yourturn}

% =====================================================================
\section*{\sffamily Ladder 4 \textbullet\ Integrated rate laws: which plot is straight?}
\ced{5.3}

Given concentration-versus-time data, the order is identified by which
plot comes out linear.

\begin{workedexample}[I do: the three tests]
State which plot is linear for each order, and what the slope gives.

\textbf{Zero order:} $[\ce{A}]$ against $t$ is linear. Slope $= -k$.
The concentration falls by equal \emph{amounts} in equal times.

\textbf{First order:} $\ln[\ce{A}]$ against $t$ is linear. Slope
$= -k$. The concentration falls by equal \emph{fractions} in equal times
--- which is what a constant half-life means.

\textbf{Second order:} $1/[\ce{A}]$ against $t$ is linear. Slope
$= +k$, positive because the reciprocal grows.

\textbf{How to use it:} plot all three from the same data. Exactly one
comes out straight, and that identifies the order. The slope then gives
$k$ directly --- and note the sign: negative for zero and first order,
positive for second.
\end{workedexample}

\begin{yourturn}[YOUR TURN 4 --- four questions]
\begin{enumerate}[leftmargin=*,itemsep=0.5em]
  \item A plot of $\ln[\ce{A}]$ against time is a straight line. What
        order? \workspace[1.4cm]
  \item A plot of $1/[\ce{A}]$ against time is straight with slope
        $+0.25$. Give the order and $k$. \workspace[1.6cm]
  \item For a zero-order reaction, what does a graph of $[\ce{A}]$ against
        $t$ look like? \workspace[1.6cm]
  \item Which order has a half-life independent of concentration?
        \workspace[1.4cm]
\end{enumerate}
\selfcheck{(a) first \quad (b) second, $k = \num{0.25}$
\si{\per\Molar\per\second} \quad (c) a straight line of negative slope
\quad (d) first}
\keyonly{\begin{apcanswer}(a) \textbf{First order}.
(b) \textbf{Second order}; the slope equals $k$, so
$k = \num{0.25}$~M\textsuperscript{-1}\,s\textsuperscript{-1}.
(c) A \textbf{straight line with negative slope}, falling at a constant
rate until the reactant is exhausted --- the rate does not depend on how
much is left.
(d) \textbf{First order.} Its half-life is $0.693/k$, containing no
concentration term at all.\end{apcanswer}}
\end{yourturn}

% =====================================================================
\section*{\sffamily Ladder 5 \textbullet\ Half-life}
\ced{5.3}

For a first-order reaction, $t_{1/2} = 0.693/k$ --- constant, whatever the
concentration. Every half-life removes half of whatever remains.

\begin{workedexample}[I do: half-life and successive halvings]
A first-order reaction has $k = \SI{0.0300}{\per\second}$ and starts at
\SI{0.800}{\Molar}. Find the half-life and the concentration after three
half-lives.

\[ t_{1/2} = \frac{0.693}{0.0300} = \mathbf{23.1~s} \]

\textbf{Three half-lives:}
$0.800 \to 0.400 \to 0.200 \to \mathbf{0.100~M}$, taking
$3 \times 23.1 = \SI{69.3}{\second}$.

\textbf{Why the half-life does not drift.} As the concentration falls, the
rate falls in exact proportion --- that is what first order means. Half of
a smaller amount takes exactly as long to disappear as half of a larger
one did.

\textbf{Contrast:} for a \emph{second}-order reaction the half-life
\emph{doubles} each time, because the rate falls faster than the
concentration does.
\end{workedexample}

\begin{yourturn}[YOUR TURN 5 --- four questions]
\begin{enumerate}[leftmargin=*,itemsep=0.5em]
  \item Half-life when $k = \SI{2.5e-4}{\per\second}$:
        \workspace[1.6cm]
  \item Fraction remaining after four half-lives: \workspace[1.4cm]
  \item Starting at \SI{0.600}{\Molar}, concentration after four
        half-lives: \workspace[1.6cm]
  \item Does a first-order half-life depend on the starting
        concentration? Explain. \workspace[1.8cm]
\end{enumerate}
\selfcheck{(a) \SI{2.8e3}{\second} \quad (b) $1/16$ \quad
(c) \SI{0.0375}{\Molar} \quad (d) no --- $t_{1/2} = 0.693/k$ only}
\keyonly{\begin{apcanswer}(a) $0.693/2.5\times10^{-4} =
\mathbf{2.8\times10^{3}}$~s (about \SI{46}{\minute}).
(b) $(1/2)^4 = \mathbf{1/16}$.
(c) $0.600/16 = \mathbf{0.0375}$~M.
(d) \textbf{No.} The expression $t_{1/2} = 0.693/k$ contains no
concentration term. Physically, the rate is proportional to concentration,
so halving a large amount and halving a small one take the same time. This
is unique to first order.\end{apcanswer}}
\end{yourturn}

% =====================================================================
\section*{\sffamily Ladder 6 \textbullet\ Elementary reactions}
\ced{5.4}

An elementary step describes an actual molecular event. For an elementary
step \emph{only}, the coefficients may be used as the orders.

\begin{workedexample}[I do: when coefficients are allowed]
For the elementary step \ce{2NO2 -> NO3 + NO}, write the rate law. Then
explain why the same cannot be done for the overall equation
\ce{2NO2 + F2 -> 2NO2F}.

\textbf{Elementary step:} the event is two \ce{NO2} molecules colliding,
so
\[ \text{rate} = k[\ce{NO2}]^2 \]
This is legitimate \emph{because} the step is elementary --- the equation
describes the collision itself.

\textbf{Overall equation:} it is a summary of the net change, not a
description of any single event. It says nothing about how many molecules
actually meet in the slow step. Writing
$\text{rate} = k[\ce{NO2}]^2[\ce{F2}]$ from it would be a guess, and
experiment shows the true law is $k[\ce{NO2}][\ce{F2}]$.

\textbf{Molecularity:} one molecule is unimolecular, two bimolecular.
Termolecular steps --- three particles meeting at once --- are very rare,
which is why proposed mechanisms almost never contain them.
\end{workedexample}

\begin{yourturn}[YOUR TURN 6 --- four questions]
\begin{enumerate}[leftmargin=*,itemsep=0.5em]
  \item Rate law for the elementary step \ce{A + B -> C}:
        \workspace[1.4cm]
  \item Rate law for the elementary step \ce{2A -> B}:
        \workspace[1.4cm]
  \item What is the molecularity of \ce{O3 -> O2 + O}?
        \workspace[1.4cm]
  \item Why are termolecular steps rare? \workspace[1.6cm]
\end{enumerate}
\selfcheck{(a) $k[\ce{A}][\ce{B}]$ \quad (b) $k[\ce{A}]^2$ \quad
(c) unimolecular \quad (d) three particles colliding simultaneously is
very improbable}
\keyonly{\begin{apcanswer}(a) $\text{rate} = k[\ce{A}][\ce{B}]$.
(b) $\text{rate} = k[\ce{A}]^2$.
(c) \textbf{Unimolecular} --- a single molecule decomposes.
(d) Because a reaction event requires the particles to be in the same
place at the same instant \emph{and} correctly oriented. Two-body
collisions are common; three particles meeting simultaneously with the
right geometry is vastly less probable, so such steps are almost never
rate-competitive.\end{apcanswer}}
\end{yourturn}

% =====================================================================
\section*{\sffamily Ladder 7 \textbullet\ The collision model}
\ced{5.5}

A collision produces reaction only if it carries at least the activation
energy \emph{and} the molecules are correctly oriented. Everything that
changes a rate changes one of those two.

\begin{workedexample}[I do: why a small warming does so much]
A rise of \SI{10}{\celsius} can roughly double a rate, though it raises
the average speed only a few percent. Explain.

\textbf{What actually matters} is the \emph{fraction} of collisions
carrying at least $E_a$, and that fraction lies far out in the
high-energy tail of the distribution of molecular energies.

\textbf{Warming shifts and broadens the whole distribution.} Near the
average that is a small change. Out in the tail, where the threshold sits,
the population rises steeply --- easily doubling.

\textbf{Collision frequency rises too, but only slightly} --- a few
percent, in line with the speed. Essentially all of the rate increase
comes from the energy distribution, not from more collisions.

\textbf{The answer that earns half marks:} ``molecules move faster and
collide more often''. True, but it is the smaller effect; the marks are
for the fraction above $E_a$.
\end{workedexample}

\begin{yourturn}[YOUR TURN 7 --- four questions]
\begin{enumerate}[leftmargin=*,itemsep=0.5em]
  \item State the two requirements for a successful collision.
        \workspace[1.4cm]
  \item Why does increasing concentration increase the rate?
        \workspace[1.6cm]
  \item Why does powdering a solid reactant speed up its reaction?
        \workspace[1.6cm]
  \item Which effect of warming matters more: more collisions, or more
        energetic ones? \workspace[1.8cm]
\end{enumerate}
\selfcheck{(a) enough energy, correct orientation \quad (b) more frequent
collisions \quad (c) more exposed surface area \quad (d) more energetic
ones, by far}
\keyonly{\begin{apcanswer}(a) Energy at least equal to the
\textbf{activation energy}, and correct \textbf{orientation}.
(b) More particles per unit volume means they meet more often, so the
collision frequency rises and more of those collisions are successful per
second.
(c) Powdering exposes far more \textbf{surface area}, so more particles
are available to collide at the solid--fluid boundary at any moment.
(d) \textbf{More energetic ones, by far.} Collision frequency rises only a
few percent with a modest warming, while the fraction of collisions
exceeding $E_a$ can double, because that fraction sits in the steep tail
of the energy distribution.\end{apcanswer}}
\end{yourturn}

% =====================================================================
\section*{\sffamily Ladder 8 \textbullet\ Reading an energy profile}
\ced{5.6} \ced{5.10}

The vertical axis is potential energy. The peak is the transition state;
the climb from reactants to peak is $E_a$; the difference between the two
ends is $\Delta H$.

\begin{workedexample}[I do: everything on one diagram]
A profile rises from reactants at \SI{50}{\kilo\joule} to a peak at
\SI{170}{\kilo\joule}, then falls to products at \SI{20}{\kilo\joule}.
Find $E_a$ forward, $E_a$ reverse and $\Delta H$, and say whether it is
exothermic.

\textbf{Forward $E_a$:} $170 - 50 = \mathbf{120~kJ}$ --- reactants up to
the peak.

\textbf{Reverse $E_a$:} $170 - 20 = \mathbf{150~kJ}$ --- products up to
the same peak.

\textbf{$\Delta H$:} $20 - 50 = \mathbf{-30~kJ}$, products minus
reactants. Negative, so \textbf{exothermic} --- and the products sit lower
than the reactants, which is the visual signature.

\textbf{The relationship worth remembering:}
$\Delta H = E_{a,\text{forward}} - E_{a,\text{reverse}} = 120 - 150 =
-30$~kJ. \checkmark

\textbf{A multi-step profile} has one peak per step, with shallow minima
between them where intermediates sit. The \emph{highest} peak is the
rate-determining step.
\end{workedexample}

\begin{yourturn}[YOUR TURN 8 --- four questions]
\begin{enumerate}[leftmargin=*,itemsep=0.5em]
  \item Reactants \SI{30}{\kilo\joule}, peak \SI{110}{\kilo\joule},
        products \SI{75}{\kilo\joule}. Forward $E_a$?
        \workspace[1.4cm]
  \item $\Delta H$ for that reaction, and is it endo- or exothermic?
        \workspace[1.6cm]
  \item On a two-step profile, how do you identify the rate-determining
        step? \workspace[1.6cm]
  \item What sits in the shallow minimum between two peaks?
        \workspace[1.4cm]
\end{enumerate}
\selfcheck{(a) \SI{80}{\kilo\joule} \quad (b) $+45$~kJ, endothermic \quad
(c) the step with the highest peak \quad (d) an intermediate}
\keyonly{\begin{apcanswer}(a) $110 - 30 = \mathbf{80}$~kJ.
(b) $75 - 30 = \mathbf{+45}$~kJ, so \textbf{endothermic} --- the products
lie above the reactants.
(c) It is the step whose peak is \textbf{highest above the preceding
minimum} --- the largest barrier to climb, and therefore the slowest step,
which sets the overall rate.
(d) An \textbf{intermediate} --- a real species, produced by one step and
consumed by a later one. Being in a minimum it is genuinely stable enough
to exist, unlike a transition state, which sits at a
\emph{maximum}.\end{apcanswer}}
\end{yourturn}

% =====================================================================
\section*{\sffamily Ladder 9 \textbullet\ Mechanisms and the rate-determining step}
\ced{5.7} \ced{5.8} \ced{5.9}

A mechanism is valid only if its steps add to the overall equation
\emph{and} its slow step predicts the observed rate law.

\begin{workedexample}[I do: testing a mechanism]
Test this mechanism for \ce{NO2 + CO -> NO + CO2}, whose observed rate law
is $k[\ce{NO2}]^2$:
\begin{align*}
\text{Step 1:}\quad & \ce{NO2 + NO2 -> NO3 + NO} \qquad \text{(slow)}\\
\text{Step 2:}\quad & \ce{NO3 + CO -> NO2 + CO2} \qquad \text{(fast)}
\end{align*}

\textbf{Do the steps add?} Summing:
\ce{2NO2 + NO3 + CO -> NO3 + NO + NO2 + CO2}. Cancel \ce{NO3} from both
sides and one \ce{NO2}:
\ce{NO2 + CO -> NO + CO2}. \checkmark

\textbf{Does the slow step give the observed law?} Step 1 is elementary
and bimolecular in \ce{NO2}, so
$\text{rate} = k[\ce{NO2}]^2$. \checkmark

\textbf{The striking prediction:} \ce{CO} does not appear in the rate law
at all, even though it is a reactant. It enters only in the \emph{fast}
step, after the bottleneck --- so adding more of it does not speed the
reaction up. That is exactly what the experiment finds, and it is strong
evidence for this mechanism.
\end{workedexample}

\begin{yourturn}[YOUR TURN 9 --- four questions]
\begin{enumerate}[leftmargin=*,itemsep=0.5em]
  \item In the mechanism above, identify the intermediate.
        \workspace[1.4cm]
  \item Which step determines the rate, and why? \workspace[1.6cm]
  \item Why does \ce{CO} not appear in the rate law? \workspace[1.6cm]
  \item What two tests must any proposed mechanism pass?
        \workspace[1.8cm]
\end{enumerate}
\selfcheck{(a) \ce{NO3} \quad (b) step 1, the slow one \quad (c) it
appears only after the slow step \quad (d) steps sum to the overall
equation; slow step gives the observed rate law}
\keyonly{\begin{apcanswer}(a) \textbf{\ce{NO3}} --- produced in step 1 and
consumed in step 2, so it appears in the mechanism but not in the overall
equation.
(b) \textbf{Step 1}, the slow step. The overall reaction cannot proceed
faster than its slowest stage, so that step is the bottleneck.
(c) Because it is consumed only in the \emph{fast} step, which happens
after the bottleneck. Whatever \ce{NO3} step 1 produces is mopped up
essentially instantly, so the supply of \ce{CO} never limits anything and
its concentration does not affect the overall rate.
(d) (1) The steps must \textbf{sum to the overall balanced equation}, with
intermediates cancelling. (2) The \textbf{slow step must predict the
experimentally observed rate law}. Failing either rules the mechanism out
--- though passing both does not prove it correct, only that it is
consistent.\end{apcanswer}}
\end{yourturn}

% =====================================================================
\section*{\sffamily Ladder 10 \textbullet\ Intermediates versus catalysts}
\ced{5.7} \ced{5.11}

Both appear inside a mechanism and neither appears in the overall
equation --- but they are opposites. An intermediate is
\emph{made then used}; a catalyst is \emph{used then remade}.

\begin{workedexample}[I do: telling them apart]
In this mechanism, identify the catalyst and the intermediate:
\begin{align*}
\text{Step 1:}\quad & \ce{X + A -> XA} \\
\text{Step 2:}\quad & \ce{XA + B -> C + X}
\end{align*}

\textbf{\ce{X} is the catalyst.} It is \emph{consumed} in step 1 and
\emph{regenerated} in step 2 --- present at the start, present at the end,
unchanged overall.

\textbf{\ce{XA} is the intermediate.} It is \emph{produced} in step 1 and
\emph{consumed} in step 2 --- absent at the start, absent at the end.

\textbf{The test is the order of events.} Consumed first, then produced:
catalyst. Produced first, then consumed: intermediate.

\textbf{Overall:} \ce{A + B -> C}, with \ce{X} appearing nowhere ---
which is why a catalyst is never written into the equation, only above the
arrow.
\end{workedexample}

\begin{yourturn}[YOUR TURN 10 --- four questions]
\begin{enumerate}[leftmargin=*,itemsep=0.5em]
  \item A species is produced in step 1 and consumed in step 3. What is
        it? \workspace[1.4cm]
  \item A species is consumed in step 1 and regenerated in step 3. What is
        it? \workspace[1.4cm]
  \item Why does neither appear in the overall equation?
        \workspace[1.6cm]
  \item Can a catalyst appear in the rate law? \workspace[1.6cm]
\end{enumerate}
\selfcheck{(a) an intermediate \quad (b) a catalyst \quad (c) both cancel
when the steps are added \quad (d) yes, if it acts at or before the slow
step}
\keyonly{\begin{apcanswer}(a) An \textbf{intermediate}.
(b) A \textbf{catalyst}.
(c) Because both appear on \emph{both} sides when the steps are summed ---
once as a product and once as a reactant --- so they cancel algebraically.
(d) \textbf{Yes.} A catalyst is not a reactant overall, but if it
participates in or before the rate-determining step, its concentration
does affect the rate and can appear in the experimental rate law. This is
common in enzyme and acid catalysis.\end{apcanswer}}
\end{yourturn}

% =====================================================================
\section*{\sffamily Ladder 11 \textbullet\ How a catalyst works}
\ced{5.11}

A catalyst provides an \emph{alternative pathway} with a lower activation
energy. It does not push the old reaction harder --- it opens a new route.

\begin{workedexample}[I do: what changes and what does not]
State precisely what a catalyst changes and what it leaves alone.

\textbf{Changes: the activation energy.} The new pathway has a lower
barrier, so at any given temperature a larger fraction of collisions
carries enough energy. The rate rises, often by orders of magnitude.

\textbf{Does \emph{not} change $\Delta H$.} Reactants and products sit at
exactly the same energies as before. Only the height of the hill between
them is reduced --- the two ends are untouched.

\textbf{Does \emph{not} change the equilibrium position or $K$.} The
forward and reverse barriers are lowered by the \emph{same} amount, so
both rates rise in the same proportion. Equilibrium arrives sooner, in
exactly the same place.

\textbf{Two answers that earn zero:} ``it lowers the energy of the
reactants'' and ``it shifts the equilibrium toward products''. Both
describe things a catalyst cannot do.
\end{workedexample}

\begin{yourturn}[YOUR TURN 11 --- four questions]
\begin{enumerate}[leftmargin=*,itemsep=0.5em]
  \item What does a catalyst do to the activation energy?
        \workspace[1.4cm]
  \item What does it do to $\Delta H$? \workspace[1.4cm]
  \item What does it do to the yield at equilibrium?
        \workspace[1.4cm]
  \item Explain why a catalyst speeds the reverse reaction too.
        \workspace[1.8cm]
\end{enumerate}
\selfcheck{(a) lowers it, via a new pathway \quad (b) nothing \quad
(c) nothing \quad (d) the same lowered barrier is crossed either way}
\keyonly{\begin{apcanswer}(a) \textbf{Lowers} it, by providing an
alternative pathway --- not by altering the original one.
(b) \textbf{Nothing.} $\Delta H$ depends only on the energies of reactants
and products, which are unchanged.
(c) \textbf{Nothing.} The equilibrium position and $K$ are unaffected;
only the time taken to reach equilibrium is shorter.
(d) Because both directions pass over the \emph{same} transition state. A
pathway with a lower barrier is a lower barrier whichever way it is
crossed, so the reverse activation energy falls by the same amount and the
reverse rate rises in the same proportion. That is precisely why the
equilibrium position cannot shift.\end{apcanswer}}
\end{yourturn}

% =====================================================================
\section*{\sffamily Ladder 12 \textbullet\ Kinetics versus thermodynamics}
\ced{5.5} \ced{5.6}

Two independent questions: \emph{does it want to happen} (thermodynamics)
and \emph{will it happen quickly} (kinetics). The answers need not agree.

\begin{workedexample}[I do: favourable but immeasurably slow]
Diamond converting to graphite has
$\Delta G^\circ = \SI{-2.9}{\kilo\joule\per\mole}$ --- thermodynamically
favourable at room temperature. Explain why diamonds do not visibly
change.

\textbf{Thermodynamics says it should happen.} A negative $\Delta G$ means
graphite is the more stable form, and the conversion would release energy.

\textbf{Kinetics says it will not, on any human timescale.} Rearranging
diamond's covalent network requires breaking very strong \ce{C-C} bonds ---
an enormous activation energy. At room temperature essentially no atom has
that much energy, so the rate is indistinguishable from zero.

\textbf{The general statement:} $\Delta G$ tells you about the
\emph{destination}; $E_a$ tells you about the \emph{journey}. A large
negative $\Delta G$ with a large $E_a$ describes a reaction that is
strongly favourable and completely stuck --- which is why fuels can sit
safely in a tank in air.
\end{workedexample}

\begin{yourturn}[YOUR TURN 12 --- four questions]
\begin{enumerate}[leftmargin=*,itemsep=0.5em]
  \item Petrol and oxygen do not react at room temperature although the
        reaction is highly favourable. Why? \workspace[1.8cm]
  \item What does a spark supply? \workspace[1.4cm]
  \item Does a large equilibrium constant mean a fast reaction?
        \workspace[1.6cm]
  \item Name two ways to speed a reaction without changing $\Delta H$.
        \workspace[1.6cm]
\end{enumerate}
\selfcheck{(a) high activation energy \quad (b) energy to surmount $E_a$
\quad (c) no \quad (d) raise the temperature; add a catalyst}
\keyonly{\begin{apcanswer}(a) The \textbf{activation energy is high}, so at
room temperature almost no collisions carry enough energy. The mixture is
under \textbf{kinetic control}: favourable but stuck.
(b) Enough \textbf{energy} to get a small number of molecules over the
activation barrier. Once started, the reaction is exothermic and supplies
the energy to keep itself going.
(c) \textbf{No.} $K$ measures how \emph{far} a reaction proceeds, not how
\emph{fast}. The two are independent: a reaction with an enormous $K$ can
be immeasurably slow if $E_a$ is large.
(d) \textbf{Raise the temperature} (increases the fraction of collisions
above $E_a$) and \textbf{add a catalyst} (lowers $E_a$ via a new
pathway). Increasing concentration or surface area also works. None of
them alters $\Delta H$.\end{apcanswer}}
\end{yourturn}

% =====================================================================
\section*{\sffamily Where you stand}

Tick a ladder only if all four YOUR TURN questions were right first time.

\renewcommand{\arraystretch}{1.30}
\noindent\begin{tabular}{@{}c p{6.0cm} c >{\raggedright\arraybackslash}p{4.9cm}@{}}
\toprule
\textbf{First try?} & \textbf{Skill} & \textbf{Ladder} &
  \textbf{If not, re-read\ldots}\\
\midrule
$\square$ & expressing rate & 1 & divide by coefficients \\
$\square$ & rate law from data & 2 & hold the others constant \\
$\square$ & $k$ and its units & 3 & units follow the order \\
$\square$ & which plot is straight & 4 & $\ln$ first, $1/[\ce{A}]$ second \\
$\square$ & half-life & 5 & first order: independent of $c$ \\
$\square$ & elementary steps & 6 & coefficients allowed only here \\
$\square$ & collision model & 7 & the tail above $E_a$ \\
$\square$ & energy profiles & 8 & peak minus start \\
$\square$ & mechanisms & 9 & sum, then slow step \\
$\square$ & intermediate vs catalyst & 10 & which comes first \\
$\square$ & how catalysts work & 11 & new path, same $\Delta H$ \\
$\square$ & kinetics vs thermodynamics & 12 & how fast vs how far \\
\bottomrule
\end{tabular}

\begin{apcnote}[Scoring yourself honestly]
12/12: move on to the unit worksheets and the free-response set.
9--11: solid --- redo the missed ladders tomorrow from a blank page.
8 or fewer: the recurring root causes here are exactly three ---
(1) reading reaction orders off the balanced equation instead of from
data or the slow step, (2) quoting a rate constant with no units, and
(3) answering a rate question with a stability argument. Fix those three
and most of these ladders fall together.
\end{apcnote}

\end{document}
